\section{The Bulk-Cusp Interface}

Spin and gauge structure live in the four-dimensional bulk coin, the $24$-cell with its $D_4$ direction set.
Matter, on the other hand, propagates on the three-dimensional cubic cusp. The natural question is how the two
are related, and whether the four-dimensional spinor structure can survive on a three-dimensional surface. The
answer is that the relation is a \emph{projection}, the standard dimensional reduction of a slice. Choosing the
cusp picks a radial direction, and everything in the bulk reduces along it in the textbook way.

\subsection{The spinor projects}

The bulk's spacetime rotation group is $SO(4) = SU(2) \times SU(2)$. Choosing a horosphere picks out a radial
direction, and that choice breaks $SO(4) \to SO(3)$. The surviving $SO(3)$ is the cusp's three-dimensional
rotation group, realized as the diagonal $SU(2)$ inside the product. Under this breaking the four-dimensional
Dirac spinor, which carries the representation $(2,1) + (1,2)$, branches to $2 + 2$. Each piece is a
three-dimensional Pauli spinor of spin one half. The branching is verified by computing the quadratic Casimir,
which reads $3/4$ on each piece, the correct value for spin one half \cite{pollard2026vibetest}.

So the four-dimensional bulk Dirac spinor projects to three-dimensional Pauli spinors. The bulk spin becomes the
cusp spin. This is exactly the textbook reduction of a $(3+1)$-dimensional Dirac fermion to its
three-dimensional spatial spinor, with $\mathrm{Spin}(3) = SU(2)$. The spinor coin and physical space are not a
mismatch to be explained away. The cusp spin is the shadow of the bulk spin, and the shadow is computed rather
than asserted.

\begin{result}[The spinor projects $4 \to 2+2$]
Under $SO(4) \to SO(3)$ from the radial choice, the bulk Dirac spinor $(2,1)+(1,2)$ branches to $2+2$, two 3D
Pauli spinors each of spin one half, with Casimir $3/4$ on each. The bulk spin projects to the cusp spin in the
standard $(3+1)D \to 3D$ reduction \cite{pollard2026vibetest}.
\end{result}

\subsection{The directions project}

The same projection acts on the direction set. The $24$ bulk directions, projected onto the cusp tangent space
orthogonal to the radial direction, collapse to a smaller three-dimensional set. The projected set holds roughly
nine to twelve distinct directions, richer than the naive cubic six \cite{pollard2026vibetest}. This projected
direction set is the cusp coin, and it coarse-grains to the $\{4,3,4\}$ cubic continuum in the infrared. A
richer projected coin carries a useful consequence. The cusp anisotropy is even smaller than the naive
cubic estimate would suggest, because the extra directions average the propagation more evenly than six axes
alone would.

\subsection{The symmetry chain and the anisotropy}

The symmetry reduces in a chain. The bulk coin carries $SO(8)$ and the bulk spacetime carries $SO(4)$. On the
cusp these reduce to the octahedral cubic point group, a finite subgroup of $SO(3)$, which then restores to full
$SO(3)$ in the infrared as the lattice coarse-grains. The exact symmetry of the cusp is the finite octahedral
group. This is the reason the cusp carries a small order-four anisotropy. The bulk, with its $D_4$ and $F_4$
structure, is isotropic to order four and protects its own directions, but that protection lives in four
dimensions. The cusp's exact symmetry is the finite octahedral group, which is anisotropic at order four, so the
order-four residual is the cusp's own and is not shielded by the bulk $F_4$.

\begin{proposition}[The interface is a projection]
The bulk-cusp relation is the dimensional reduction of a slice along the radial direction. Spin reduces
$SO(4) \to SO(3)$ with $4 \to 2+2$, the $24$ bulk directions project to the $\sim 9$ to $12$ direction cusp
coin, and full $SO(3)$ is restored in the infrared. The internal $SO(8)$ structure rides along as internal
quantum numbers.
\end{proposition}


The interface is not an unbridged gap between a four-dimensional coin and a three-dimensional space. It is a
projection, made precise. Spin reduces $SO(4) \to SO(3)$, the directions reduce to the cusp coin, and full
$SO(3)$ is restored in the infrared. The internal $SO(8)$ structure, which carries the gauge content and the
generation structure, rides along as internal quantum numbers and is untouched by the spatial projection. What
remains is bookkeeping, the exact lattice-level coin and the detailed carry-over of the $SO(8)$ labels, rather
than a missing mechanism. The mechanism is the standard reduction of a slice, and it does what it needs to do.
